\section{Historiography of the Right}

\begin{quotex}
This essay by Julius Evola was originally published in the journal ``Roma'' on 7 August 1973 under the title ``Storiografia di Destra''.
\end{quotex}

In developing some considerations on the European meaning that can be attributed to Donoso Cortes, the Spanish thinker and an interesting type of the political man, who developed his activity in the period of the first European revolutionary and socialist movements, a noted German historian, Carl Schmitt, pointed out that while the left has systematically elaborated and perfected their historiography as the general background for their destructive action, nothing similar happened in the opposed camp of the Right. There the whole is reduced to some sporadic example in no comparable way, through consistency, radicalism, and breadth of horizons, to that which Marxism and the left possessed for some time.

To a large measure, this is correct. Actually, the only noted history that had more influence, excluding that of Marxist intonation, is essentially of liberal, illuminist, and masonic origins. It relates to those ideologies of the Third Estate that served only to prepare the ground for radical movements of the left, having itself an essentially anti-traditional foundation. An historiography of the Right still waits to be written, and that constitutes one of our signs of inferiority in respect to the ideologies and actions of unrest of the left. In particular, the so-called ``homeland history'' \emph{storica patria} in vogue cannot fill the lacuna, because apart from certain of its possible national colorations and the moving re-evocations of events and heroic figures, it itself is affected, in a large measure, by suggestions of thought that is not of the true Right and, especially, because it cannot stand in comparison, as far as the breadth of horizons, to the historiography of the left.

This is the fundamental point. In fact, we must recognize that the historiography of the left has been able to take the view on the essential dimensions of history: beyond conflicts and episodic political developments, beyond the history of nations it has been able to see the general process and realized fundamentals in recent centuries, in the sense of the transition from one type of civilization and society to another. That the basis of the interpretation was, in such regard, economic and classist, does not take away any of the breadth of the description of the whole treated by that historiography. That, as the essential reality beyond the contingent and particular, indicates to us, in the course of history, the end of the feudal and aristocratic civilization, the arrival of the liberal, capitalist, and industrial bourgeois civilization and, after this, the announcement and incipient realization of a socialist, Marxist, and finally communist civilization. Here the revolutions of the Third Estate and the Fourth Estate are recognized in their natural causal and tactical concatenation. The idea of higher level processes which, without intending or knowing it, served the more or less ``sacred'' egoisms of the people, the rivalries and ambitions of those who believed they ``make history'' without leaving the field of the particular, such ideas are certainly considered. They study exactly the transformations of the whole of the social structure and civilization that are the direct effect of the play of historical forces, rightly relegating the history of nations to the simple ``bourgeois'' phase of general development. (In fact, the ``nations'' did not arise in history as its subject but started from the revolution of the Third Estate and as its consequence.)

Measured by the historiography of the left, which is characteristic of other tendencies, it therefore appears superficial, episodic, bi-dimensional, sometimes even frivolous. A historiography of the Right should embrace the same horizons as Marxist historiography, with the will to gather the real and essential elements of the historical process developed in recent centuries, beyond myths, superstructures, and also ordinary news. This, naturally, by inverting the signs and the perspective, i.e., seeing in the essential and convergent processes of recent history not the phases of a political and social progress but rather that of a general subversion. As is logical, the economic-material premise also needs to be eliminated, recognizing \emph{homo oeconomicus} and the presumed inexorable determinism of the various means of production as mere fictions.

Much vaster, deeper, and complex forces were and are in the action of history. The particulars and even the myth of so-called ``primordial communism'' must be rejected and must be countered by the idea of organizations based predominately on a principle of a pure, sacred, and traditional spiritual authority as the civilizations that preceded those of the feudal aristocratic type. But, apart from this, we repeat, an historiography of the Right will recognize, not less than that of the left, the succession and concatenation of general distinct supra-national phases, which regressively led up to disorder and the current subversions: and this will be, for it, the basis for the interpretation of individual facts and upheavals, always attentive to the effect it produced in the total picture.

Here it is impossible to indicate, not even with any examples, all the fecundity of such a method, the unsuspected light that it would throw on a large number of events. The political-religious conflicts of the imperial Middle Ages, the constant schismatic action of France, the relations between England and Europe, the true meaning of the ``achievements'' of the French Revolution and little by little up to the episodes that particularly interest us, as the effective turn of the revolt of the Communes, the double aspect of the Risorgimento which national movement, triggered however by ideologies of the Third Estate, the significance of the Holy Alliance and the efforts of Metternich (that last great European), then that of the first world war with the ricochet action of its ideologies, the discrimination of the positive and negative in the national revolutions that recently have been asserted in Italy and in Germany, and so on, in order to finally reach a vision in conformance with the naked reality of the true forces in battle today for the control of the world. Here is a choice of suggestive arguments, among so many, to which the historiography of the Right could be applied, revolutionizing the views that most people are accustomed to have on everything through the effect of historiography of the opposite orientation, and acting in an illuminating way.

An historiography so formulated, looking at the universal, would then be especially at the height of the times if it is true that, through the effect of objective irreversible processes, alignments increasingly stand out today that are not just ethnic and particular political and closed unities. Except that unfortunately only from the hoped for historiography would an increase in consciousness come. Its effectiveness as well as practice could have with difficulty, in the current state of things, the goal of an action, of a global and inexorable battle against the forces that stand for sweeping away what little that still remains of true European tradition. In fact, it could so happen that, as an opposing party, an international of the Right could exist, organized and equipped with strength like the communists. Now, unfortunately we know that through the lack of men of high stature and sufficient authority, for the predominance of party interests and small ambitions, for the lack of true principles and not through the lack of intellectual courage, a unitary alignment of the Right so far has not been possible to be constituted not even in our Italy and that only recently have initiatives of this type been announced.


\flrightit{Posted on 2012-11-05 by Cologero}

\begin{center}* * *\end{center}

\begin{footnotesize}
\begin{sffamily}
\texttt{Cologero on 2012-11-05 at 21:34 said:}

Since this essay was written, there has been no progress on this proposed task. The historiography of the left is accepted as obviously true: to wit, history, or at least modern history, has been the progressive growth of knowledge, freedom, and tolerance; it has overcome a series of enslavement, oppression, and discrimination.

Opposed to this, Evola expresses the counter-view of Tradition. Modern history has been the decline from a spiritual, heroic, hierarchical civilization to a materialist, undifferentiated, meaningless existence. The cause is the ``degeneration of castes''.

He assumes that current events will be understood in this light, at least by the few who are up to this task. Are there any takers?

Nowhere, even with the micro-surge in the interest in Tradition in some quarters, is such an historiography being written. I assume the global right that Evola referred to was GRECE, or the New Right as it is more commonly called. Unfortunately, the perspective of this movement is far from what Evola is describing here.

\hfill

\texttt{Mihai on 2012-11-06 at 04:47 said:}

Here is part of a comment I wrote elsewhere concerning the new right:

``The new-right claims to abhor modernity's reduction of everything to quantity and to promote quality and principles. Yet, their methods are as quantitative as those of the left and they believe that the key to victory is to ``convince'' as many people as possible and to have greater numbers than their leftist opponents. It seems that, for the new-right, the way to quality is to have the largest quantity on its side. Perplexing.

Last, but not least: The new-right claims to stand for solid principles and a superior vision of life. Yet, reading the likes of Alain de Benoist one only finds there a curious mixture of ``rightist'' jargon, a few Nietzschean slogans perpetuated mechanically, some vulgarized ideas from Evola and loads of mystifications.

Worst of all, the new-right is greatly infiltrated by all sorts of pathetic neo-paganisms, crowleyan occultisms, new-agey type of naturalism and, worst of all, darwinian evolutionism, this last being the main pillar on which the modern mentality is built.

Now I do not doubt that this right that is at stake here may, in the future, gain some political support and even find itself in a position of power. Fortunes are always changing. But then what ? Are we to replace the current lie with some even more dangerous half-truths? Because this is what the so-called ``principles'' of the new-right amount to. And then we'll replace the caprice of the mass and of the lowest common denominator with the caprice of a single individual or that of a small circle of a pseudo-elite. Nazism was something like this, to be honest.

\hfill

\texttt{Graham on 2012-11-06 at 11:41 said:}

Plinio Correa de Oliveira's `Revolution and Counter-Revolution' could serve, at least, as a starting point for such a project.

\hfill

\texttt{aperion on 2012-11-10 at 18:40 said:}

GRECE is old news, the manifestations of the ``New Right'' (both in Europe and in the U.S) have come a long way and are open to ideas aligning with Tradition. what's more is that there is an audience for this material. What is lacking is leadership, outreach and the willingness to get knee deep in the mud by engaging with people both real and online. Evola hints at the prospects of networking with sympathizers in MATR (or what was politically possible at that time), as it is true that men are more often won over by example then sterile dialectical intercourse.

Now the time has come for Cologero to begin such a historiography of the Right, in a publishable print format.

\hfill

\texttt{Ulysses on 2016-11-06 at 15:03 said:}

``Are there any takers?''

I had been trying to write a book titled ``Catholicism and the Causation of the Modern Deviation'', but I didn't have the necessary knowledge. By the time I do insha'Allah, I may have a new title, like ``Genghis Khan and the Making of the Modern World''.
\url{https://www.amazon.com/Genghis-Khan-Making-Modern-World/dp/B005D60UOC}


\hfill

\texttt{Cologero on 2019-06-26 at 19:53 said:}

That's not at all on the ``right track'' ... there is nothing about the ``third dimension'' of history.

Stay away from Bolton, or else stay away from here.


\hfill

\end{sffamily}
\end{footnotesize}

